\section{Цель работы}

Ознакомиться с основными приемами работы с
документацией при составлении программ для микроконтроллеров

\section{Постановка задачи}
Создать новый проект в Keil μVision5 и разработать программу для микроконтроллера STM32F200, которая включает и выключает светодиод, а также вторую часть задания с “бегущим светодиодом”, а именно последовательного включения всех светодиодов.

\section{Алгоритм работы программы}

Данный раздел описывает последовательность действий, необходимых для разработки и реализации проекта на базе микроконтроллера STM32F207IGHx в среде Keil~µVision5. Алгоритм включает этапы инициализации проекта, настройки портов ввода–вывода, реализации мигания светодиода и модификации программы для создания эффекта «бегущего огня».

\subsection{Инициализация проекта}

\begin{enumerate}
  \item Создать новый проект в среде Keil~µVision5.
  \item Выбрать микроконтроллер семейства STM32F207IGHx.
  \item Выполнить базовую конфигурацию проекта:
  \begin{itemize}
      \item подключить поддержку ядра CMSIS (CMSIS~CORE);
      \item добавить стартовый файл;
      \item подключить необходимые модули HAL.
  \end{itemize}
\end{enumerate}

\subsection{Настройка портов GPIO}

Для работы со светодиодами требуется настроить порты \texttt{GPIOG}, \texttt{GPIOH} и \texttt{GPIOI}.

\begin{enumerate}
  \item Разрешить тактирование соответствующего порта, установив нужные биты в регистре \texttt{RCC\_AHB1ENR}.
  \item Настроить режим работы выводов как «выход» через регистр \texttt{GPIOx\_MODER}.
  \item Определить параметры выхода:
  \begin{itemize}
      \item тип выхода — push--pull (\texttt{GPIOx\_OTYPER});
      \item низкая скорость переключения (\texttt{GPIOx\_OSPEEDR}).
  \end{itemize}
\end{enumerate}

\subsection{Реализация мигания светодиодов}

\begin{enumerate}
  \item Включить тактирование портов ввода-вывода (GPIOG, GPIOH, GPIOI), установив соответствующие биты в регистре управления тактированием (по адресу \texttt{0x40023830}).
  \item Настроить используемые пины (PG6, PG7, PG8, PH2, PH3, PH6, PH7, PI10) на режим вывода (general purpose output mode), записав значение \texttt{01} в соответствующие биты регистров \texttt{MODER} каждого порта.
  \item Установить бит нужного пина (начиная с PH3) в регистре выходных данных (\texttt{ODR}) в состояние «1», включая светодиод.
  \item Выполнить пустой цикл (\texttt{for}) для создания программной задержки, необходимой для визуального восприятия.
  \item Сбросить бит этого же пина в регистре \texttt{ODR} в «0», выключив светодиод.
  \item Повторить шаги 3--5 для всех остальных задействованных пинов в заданной последовательности для получения эффекта поочередного мигания.
\end{enumerate}


\newpage

\section*{Код программы}

\begin{lstlisting}[style=cstyle, caption={Основная функция программы}]
int main()
{
 int i;
 unsigned long int j;
 i = 0;
 j = 0;
 
 // 1. Включение тактирования портов ввода-вывода (регистр RCC_AHB1ENR, адрес 0x40023830)
 // Чтобы включить тактирование порта, нужно установить соответствующий бит в 1.
 // Вычисление значения: 2 в степени (номер бита). Перевод в 16-рич систему:
 // Бит 6 (порт G) = 2^6 = 64 = 0x40
 // Бит 7 (порт H) = 2^7 = 128 = 0x80
 // Бит 8 (порт I) = 2^8 = 256 = 0x100
 *(unsigned long*)(0x40023830) |= 0x40;  // Включаем тактирование GPIOG
 *(unsigned long*)(0x40023830) |= 0x80;  // Включаем тактирование GPIOH
 *(unsigned long*)(0x40023830) |= 0x100; // Включаем тактирование GPIOI
 
 // Программная задержка для стабилизации частоты перед настройкой портов
 for (i = 0; i < 4; i++) {}

 // 2. Настройка направления работы пинов (регистры MODER)
 // В регистре MODER на каждый пин выделяется по 2 бита (с номерами 2*N и 2*N + 1, где N - номер пина).
 // Для режима выхода (Output) нужно записать комбинацию "01", то есть:
 // - сбросить старший бит (2*N + 1) в 0 (через операцию И с инвертированной маской)
 // - установить младший бит (2*N) в 1 (через операцию ИЛИ)
 // Пример расчета для пина 6: нужные биты 13 и 12. 
 // 2^13 = 8192 = 0x2000 (маска для сброса). 2^12 = 4096 = 0x1000 (маска для установки).

 // Порт G (базовый адрес 0x40021800)
 *(unsigned long*)(0x40021800) = (*(unsigned long*)(0x40021800) & (~0x00002000)) | (0x00001000); 
 // PG6: сбрасываем бит 13 (0x2000) и 12 (0x1000)
 
 *(unsigned long*)(0x40021800) = (*(unsigned long*)(0x40021800) & (~0x00008000)) | (0x00004000); 
 // PG7: биты 15 и 14. 2^15 = 0x8000 (сброс), 2^14 = 0x4000 (установка)
 
 *(unsigned long*)(0x40021800) = (*(unsigned long*)(0x40021800) & (~0x00020000)) | (0x00010000); 
 // PG8: биты 17 и 16. 2^17 = 0x20000 (сброс), 2^16 = 0x10000 (установка)
    
 // Порт H (базовый адрес 0x40021C00)
 *(unsigned long*)(0x40021C00) = (*(unsigned long*)(0x40021C00) & (~0x00000020)) | (0x00000010); 
 // PH2: биты 5 и 4. 2^5 = 0x20 (сброс), 2^4 = 0x10 (установка)
 
 *(unsigned long*)(0x40021C00) = (*(unsigned long*)(0x40021C00) & (~0x00000080)) | (0x00000040); 
 // PH3: биты 7 и 6. 2^7 = 0x80 (сброс), 2^6 = 0x40 (установка)
 
 *(unsigned long*)(0x40021C00) = (*(unsigned long*)(0x40021C00) & (~0x00002000)) | (0x00001000); 
 // PH6: биты 13 и 12. 2^13 = 0x2000 (сброс), 2^12 = 0x1000 (установка)
 
 *(unsigned long*)(0x40021C00) = (*(unsigned long*)(0x40021C00) & (~0x00008000)) | (0x00004000); 
 // PH7: биты 15 и 14. 2^15 = 0x8000 (сброс), 2^14 = 0x4000 (установка)

 // Порт I (базовый адрес 0x40022000)
 *(unsigned long*)(0x40022000) = (*(unsigned long*)(0x40022000) & (~0x00200000)) | (0x00100000); 
 // PI10: биты 21 и 20. 2^21 = 0x200000 (сброс), 2^20 = 0x100000 (установка)

 // 3. Управление светодиодами через регистры ODR (смещение +0x14 от базы порта)
 // Здесь каждый пин занимает ровно 1 бит (совпадает с номером пина N).
 // Значение маски рассчитывается как 2^N. Установка бита (|=) включает диод, сброс (&= ~) выключает.
	while(1) {
    // PH3: бит 3 -> 2^3 = 8 = 0x08
    *(unsigned long*)(0x40021C14) |= 0x08;   // Включение
    for (j = 0; j < 1000000; j++) {}         // Задержка для видимости эффекта
    *(unsigned long*)(0x40021C14) &= ~0x08;  // Выключение

    // PH6: бит 6 -> 2^6 = 64 = 0x40
    *(unsigned long*)(0x40021C14) |= 0x40;   // Включение
    for (j = 0; j < 1000000; j++) {} // Задержка для видимости эффекта
    *(unsigned long*)(0x40021C14) &= ~0x40;  // Выключение

    // PH7: бит 7 -> 2^7 = 128 = 0x80
    *(unsigned long*)(0x40021C14) |= 0x80;  // Включение
    for (j = 0; j < 1000000; j++) {} // Задержка для видимости эффекта
    *(unsigned long*)(0x40021C14) &= ~0x80;  // Выключение

    // PI10: бит 10 -> 2^10 = 1024 = 0x400
    *(unsigned long*)(0x40022014) |= 0x400;  // Включение
    for (j = 0; j < 1000000; j++) {} // Задержка для видимости эффекта
    *(unsigned long*)(0x40022014) &= ~0x400;  // Выключение

    // PG6: бит 6 -> 2^6 = 64 = 0x40
    *(unsigned long*)(0x40021814) |= 0x40; // Включение   
    for (j = 0; j < 1000000; j++) {} // Задержка для видимости эффекта
    *(unsigned long*)(0x40021814) &= ~0x40;  // Выключение

    // PG7: бит 7 -> 2^7 = 128 = 0x80
    *(unsigned long*)(0x40021814) |= 0x80;   // Включение
    for (j = 0; j < 1000000; j++) {} // Задержка для видимости эффекта
    *(unsigned long*)(0x40021814) &= ~0x80;  // Выключение

    // PG8: бит 8 -> 2^8 = 256 = 0x100
    *(unsigned long*)(0x40021814) |= 0x100;  // Включение
    for (j = 0; j < 1000000; j++) {} // Задержка для видимости эффекта
    *(unsigned long*)(0x40021814) &= ~0x100; // Выключение

    // PH2: бит 2 -> 2^2 = 4 = 0x04
    *(unsigned long*)(0x40021C14) |= 0x04;   // Включение
    for (j = 0; j < 1000000; j++) {} // Задержка для видимости эффекта
    *(unsigned long*)(0x40021C14) &= ~0x04;  // Выключение

    // PG8: снова бит 8 -> 0x100
    *(unsigned long*)(0x40021814) |= 0x100;  // Включение
    for (j = 0; j < 1000000; j++) {} // Задержка для видимости эффекта
    *(unsigned long*)(0x40021814) &= ~0x100; // Выключение
	}
}
\end{lstlisting}

\newpage